% Gauss, Werke X/1 (1917), native115, printed108.
\SourceLeaf{115}\ArithmeticHead{108}{L}{ALGEBRA.\quad BRIEFWECHSEL.}
\NumberedItem{4}
\begin{center}\large\textsc{Gauss} an \textsc{Schumacher}.\end{center}
\begin{flushright}Göttingen, 1840 Juni 20.\end{flushright}
\noindent\hspace*{\parindent}Beigehend, mein theuerster Freund, schicke ich Ihnen den
\textsc{Clausen}schen Aufsatz zurück. Ich habe ihn mit Vergnügen gelesen, wenn gleich derselbe,
seinem wesentlichen Inhalt nach nur eine analytische Einkleidung desjenigen Princips ist,
welches ich im letzten (24.) Artikel meiner Schrift von 1799[*)] in geometrischer Form
angedeutet habe, sowie ich auch dort (Ende des 23. Art.) bemerkt habe, dass der Nerv davon
eigentlich mit dem \textsc{d’Alembert}ischen Princip coincidirt. Indessen macht die Tournure,
die Hr. \textsc{Clausen} ihm gegeben hat, ihn der Aufnahme in die A. N. zweifelsohne vollkommen
würdig[**)]. Ganz befriedigend und dem Rigor antiquus genügend ist übrigens die Ausführung
nicht, indem daraus allein, dass eine Grösse [sich] unaufhörlich in einerlei Sinn ändert,
z. B. wie bei \textsc{Clausen} eine reale negative, die immer noch weiter (absolut) abnehmen
kann, noch nicht evident ist, dass sie jeden auf dem Wege liegenden Werth wirklich erreichen
kann, also im \textsc{Clausen}schen Falle den 0 Werth erreichen. Ich habe diesen Umstand in
den letzten Zeilen jener Schrift selbst angemerkt, so wie zugleich, dass er sich allerdings
heben lässt, aber gerade weil diess, \textls[120]{im Geiste des Rigor antiquus}, nicht ohne
einige Umständlichkeit geschehen kann, habe ich später dieses Verfahren nicht selbst
ausgeführt.\quad . . . . . . . . . . . . . . . . . . .
\par\smallskip\noindent\rule{.25\linewidth}{.4pt}\par\smallskip
\begingroup\small
\noindent\hspace*{\parindent}[*) Werke III, S. 29.]\par
\noindent\hspace*{\parindent}[**) Der vom 16. Juni 1840 datierte Aufsatz von \textsc{Thomas Clausen},
\textit{Beweis, dass die algebraischen Gleichungen Wurzeln von der Form $a+bi$ haben,} ist
abgedruckt in den Astronomischen Nachrichten \textbf{17} (1840), Spalte 325 ff.]\par
\endgroup
\clearpage
